\begin{helper}
Fix \(I\subseteq[n]\), integers \(k,\ell_R,\ell_B\), and parameters
\(m,p,\varepsilon\).  On the projection-size event
\[
S=\{\omega:
\noderef{TwoBiteProjectionSizeEvent}(I,k,\ell_R,\ell_B,\omega)\}
\]
there is a finite disjoint partition by pre-terminal exposure histories.  For
each history \(h\) there is a representative sample \(\omega_h\) and a
recorded coordinate set \(R_h\) such that
\[
\omega\in C_h
\quad\Longleftrightarrow\quad
\omega\hbox{ has the same injection as }\omega_h
\hbox{ and agrees with }\omega_h\hbox{ on every coordinate of }R_h.
\tag{1}
\]
The recorded set \(R_h\) is the exact set of red and blue base-edge
coordinates queried before the terminal coordinate scan begins.  Equivalently,
for every \(\omega\in C_h\),
\[
e\in R_h
\quad\Longleftrightarrow\quad
e\in\noderef{TwoBitePreTerminalRecordedPairs}(\omega,\varepsilon,I).
\tag{1'}
\]
Agreement on the injection and on \(R_h\) replays every
\(\noderef{TwoBiteStageRevealedVertex}\) predicate and hence every
\(\noderef{TwoBiteProjectionPairTouched}\) predicate.  Finally, if
\(\omega\in C_h\) and
\[
e\in\noderef{TwoBiteStagedOpenPairs}(\omega,\varepsilon,I),
\]
then \(e\) is an increasing non-loop terminal coordinate in
\(\noderef{TwoBiteTerminalCoordinateUniverse}(m)\) and \(e\notin R_h\).
\end{helper}
\begin{proof}
Fix deterministic orders of the finite base-vertex set
\(\noderef{TwoBiteBaseVertex}(m)\), of ordered pairs of base vertices, and of
the terminal coordinate universe
\(\noderef{TwoBiteTerminalCoordinateUniverse}(m)\).  When two same-color base
vertices \(x,y\) are distinct, write \([x,y]\) for the corresponding tagged
coordinate with the two endpoints put in increasing order; if the colors are
different or \(x=y\), there is no such terminal coordinate.

We define a canonical pre-terminal replay for a sample \(\omega\).  The replay
keeps a finite recorded set \(R\) and, for every \(e\in R\), the truth value
\(\noderef{TwoBiteEdgeCoordinateValue}(\omega,e)\).  The order of the replay
is fixed once and for all as follows.

First, the replay reads only the injection \(\omega.2.2\).  From the injection
it computes \(\pi_R^\omega\), \(\pi_B^\omega\), every
\(\noderef{TwoBiteProjectionContains}(\omega,I,x)\), and every fiber
\(\noderef{TwoBiteBaseFiber}(\omega,x)\cap I\).  Hence it computes
\(\noderef{TwoBiteStageF1}(\omega,I,x)\) for every \(x\), since
\(\noderef{TwoBiteStageF1}\) is precisely projection containment together with
the numerical fiber-size test.

Second, the replay computes \(F_2\) by scanning the ordered pairs \((x,y)\).
If \(x\) is not projection-contained, then
\(\noderef{TwoBiteStageF2}(\omega,\varepsilon,I,x)\) is false by definition,
so no adjacency value is needed for this \(x\).  If \(x\) is
projection-contained and \(y\notin F_1\), then \(y\) cannot contribute to the
defining \(F_2\)-count.  If \(x\) and \(y\) have different colors, or if
\(x=y\), the value of \(\noderef{TwoBiteBaseAdj}(\omega,x,y)\) is false by
the definition of \(\noderef{TwoBiteBaseAdj}\) and by simplicity of the base
graphs.  The only coordinates read in this pass are therefore the distinct
same-color pairs with \(x\) projection-contained and \(y\in F_1\).  For each
such pair, before reading the edge value, the replay inserts \([x,y]\) into
\(R\), unless it was already inserted, and then reads
\(\noderef{TwoBiteEdgeCoordinateValue}(\omega,[x,y])\).

This insertion is non-circular.  At this point \(F_1\) and projection
containment have already been computed from the injection.  Since
\(y\in F_1\), the definition of
\(\noderef{TwoBiteStageRevealedVertex}\) already implies that \(y\) is
staged-revealed, independently of the yet-to-be-computed \(F_2\) table.  Thus
\([x,y]\) satisfies the incident-coordinate branch of
\(\noderef{TwoBitePreTerminalRecordedPairs}\): one endpoint, namely \(y\), is
staged-revealed and the other endpoint, namely \(x\), is
projection-contained.  The early recording of \([x,y]\) is exactly the
paper's pre-terminal query of an adjacency to an \(F_1\)-vertex; later
incident-coordinate recording may encounter the same coordinate again but
does not change \(R\).  After the ordered scan over \(y\), the replay knows
the finite set
\[
\{y:\noderef{TwoBiteStageF1}(\omega,I,y)
       \text{ and }\noderef{TwoBiteBaseAdj}(\omega,x,y)\}
\]
for every projection-contained \(x\), and therefore it knows
\(\noderef{TwoBiteStageF2}(\omega,\varepsilon,I,x)\) for every \(x\).

Third, the replay computes the staged-revealed predicate
\[
\noderef{TwoBiteStageRevealedVertex}(\omega,\varepsilon,I,x)
\Longleftrightarrow
\neg\noderef{TwoBiteProjectionContains}(\omega,I,x)
\vee\noderef{TwoBiteStageF1}(\omega,I,x)
\vee\noderef{TwoBiteStageF2}(\omega,\varepsilon,I,x).
\tag{2}
\]
Fourth, for every staged-revealed vertex \(x\), the paper reveals all
neighbors and non-neighbors of \(x\) inside the appropriate projection of
\(I\).  In tablet terms, for every same-color projection-contained endpoint
\(z\) with \(z\ne x\), the replay inserts the incident coordinate \([x,z]\)
and stores its edge value.  If \(z=x\), the relevant same-color adjacency is a
loop and is false in the simple base graph, so no terminal coordinate exists
or is queried.  The inserted coordinates are exactly the incident-coordinate
alternatives in \(\noderef{TwoBitePreTerminalRecordedPairs}\).

Fifth, with those incident values known, \(X_x(I)\) is determined for each
staged-revealed \(x\).  Indeed, for \(u\in I\), membership
\[
u\in\noderef{TwoBiteX}(\omega,I,x)
\]
means, by the definition of \(\noderef{TwoBiteX}\) and
\(\noderef{TwoBiteLiftedNeighborhood}\), that \(u\in I\) and the same-color
base edge from \(x\) to the projection of \(u\) is present.  The projection of
each \(u\in I\) is projection-contained.  If this projected endpoint is
distinct from \(x\) and has the same color, the required incident coordinate
was recorded in the preceding paragraph; if it is a loop or has the opposite
color, the relevant base adjacency is false by definition.  The replay can
therefore compute the red and blue projected images of \(X_x(I)\), then scan
the terminal coordinate universe and record exactly the red coordinates
\[
q\in\noderef{TwoBiteBasePairs}
\bigl(\pi_R^\omega(\noderef{TwoBiteX}(\omega,I,x))\bigr)
\]
and the blue coordinates
\[
q\in\noderef{TwoBiteBasePairs}
\bigl(\pi_B^\omega(\noderef{TwoBiteX}(\omega,I,x))\bigr)
\]
for staged-revealed \(x\).  These are exactly the \(X_x(I)\)-pair alternatives
in \(\noderef{TwoBitePreTerminalRecordedPairs}\).

Consequently, for every sample \(\omega\), the final output \(R(\omega)\) of
this replay satisfies
\[
e\in R(\omega)
\quad\Longleftrightarrow\quad
e\in\noderef{TwoBitePreTerminalRecordedPairs}(\omega,\varepsilon,I).
\tag{3}
\]
The forward implication follows from the two ways the replay records a
coordinate: incident coordinates satisfy the first branch of
\(\noderef{TwoBitePreTerminalRecordedPairs}\), and \(X_x(I)\)-pair coordinates
satisfy the second branch.  The reverse implication follows by the same case
split.  An incident coordinate in the first branch is recorded when the replay
handles its staged-revealed endpoint.  A coordinate in the second branch is
recorded when the replay scans pairs in \(X_x(I)\) for the witnessing
staged-revealed vertex \(x\).

Define the transcript of \(\omega\) to be
\[
\tau(\omega)=
\bigl(\omega.2.2,\ R(\omega),\
(\noderef{TwoBiteEdgeCoordinateValue}(\omega,e))_{e\in R(\omega)}\bigr).
\tag{4}
\]
Only finitely many transcripts occur: there are finitely many injections,
finitely many subsets of
\(\noderef{TwoBiteTerminalCoordinateUniverse}(m)\), and finitely many truth
tables on each subset.  Let \(H\) be the finite set of transcripts realized by
samples in \(S\).  For each \(h\in H\), choose a representative
\(\omega_h\in S\) with \(\tau(\omega_h)=h\), set \(R_h=R(\omega_h)\), and
define
\[
C_h=\{\omega:\omega.2.2=\omega_h.2.2
\text{ and }
\noderef{TwoBiteEdgeCoordinateValue}(\omega,e)
\Longleftrightarrow
\noderef{TwoBiteEdgeCoordinateValue}(\omega_h,e)
\text{ for every }e\in R_h\}.
\tag{5}
\]
This is the cylinder description (1).

We next record the deterministic uniqueness property of the replay.  Suppose
two samples \(\omega_0,\omega_1\) have the same injection and that
\(\omega_1\) agrees with \(\omega_0\) on every coordinate of \(R(\omega_0)\),
with the same truth value
\(\noderef{TwoBiteEdgeCoordinateValue}(-,e)\).  Running the replay on
\(\omega_1\) makes the same decisions as running it on \(\omega_0\), outputs
the same recorded set, and records the same truth table.  The proof is an
induction over the ordered replay steps just defined.

The injection step is identical by assumption.  Projection containment,
fibers, and \(F_1\) are functions only of that common injection, so the two
runs have the same projection-contained vertices and the same \(F_1\)-table.
During the \(F_2\)-pass, the two runs therefore consider the same ordered
candidate pairs.  Opposite-color and loop candidates are false in both runs
without reading a coordinate.  For a distinct same-color candidate with
\(x\) projection-contained and \(y\in F_1\), the first run inserts \([x,y]\).
As explained above, \(y\in F_1\) makes \(y\) staged-revealed by (2), while
\(x\) is projection-contained; hence \([x,y]\) satisfies the
incident-coordinate branch of
\(\noderef{TwoBitePreTerminalRecordedPairs}(\omega_0,\varepsilon,I)\).  By
(3), this coordinate is in \(R(\omega_0)\).  The assumed agreement on
\(R(\omega_0)\) therefore gives the same edge value for \([x,y]\) in the two
runs.  Thus every \(F_2\)-adjacency contribution is the same, and the two
\(F_2\)-counts agree.  This argument uses only the already-known \(F_1\) table
to place the queried coordinate in \(R(\omega_0)\); it does not assume the
\(F_2\)-status of \(x\) in advance.

With projection containment, \(F_1\), and \(F_2\) agreeing, equation (2) gives
the same staged-revealed vertices in the two runs.  The incident-coordinate
completion step is then the same: it scans the same staged-revealed vertices
and the same projection-contained same-color endpoints, and every coordinate
it reads is an incident branch coordinate for \(\omega_0\), hence lies in
\(R(\omega_0)\) by (3), so the stored truth values agree.  Finally consider
the \(X_x(I)\)-pair step.  For a staged-revealed \(x\) and \(u\in I\), the
projection of \(u\) is determined by the common injection.  If the relevant
base adjacency is a non-loop same-color adjacency, its coordinate was inserted
in the incident-coordinate completion step for \(\omega_0\), and therefore
belongs to \(R(\omega_0)\); agreement on \(R(\omega_0)\) gives the same
adjacency value.  Loop and opposite-color cases are false in both runs by the
definitions.  Hence the two runs compute the same sets
\(\noderef{TwoBiteX}(-,I,x)\), the same red and blue projected images, and the
same terminal \(X_x(I)\)-pair coordinates.  The replay therefore has a unique
output determined by the common injection and the truth table on the first
run's recorded set.

The cells cover \(S\).  If \(\omega\in S\), then
\(\tau(\omega)\in H\), and the representative of that transcript has the same
injection, the same recorded set, and the same recorded truth table as
\(\omega\); hence \(\omega\in C_{\tau(\omega)}\).  Conversely, if
\(\omega\in C_h\), then \(\omega\) has the same injection as
\(\omega_h\in S\).  The event
\(\noderef{TwoBiteProjectionSizeEvent}(I,k,\ell_R,\ell_B,-)\) depends only on
\(|I|\) and on the two projection images of \(I\), and those are fixed by the
injection.  Therefore \(\omega\in S\).  This proves both the covering
equivalence and the containment of each cell in \(S\).

The cells are disjoint.  If \(\omega\in C_h\cap C_{h'}\), then
\(\omega_h\), \(\omega\), and \(\omega_{h'}\) all have the same injection.
The sample \(\omega\) agrees with \(\omega_h\) on \(R_h\) and with
\(\omega_{h'}\) on \(R_{h'}\).  Since \(R_h=R(\omega_h)\), deterministic
replay uniqueness applied first to \((\omega_h,\omega)\) says that the replay
of \(\omega\) has the same recorded set and the same recorded truth table as
the replay of \(\omega_h\); therefore \(\tau(\omega)=h\).  Applying the same
argument to \((\omega_{h'},\omega)\) gives \(\tau(\omega)=h'\).  A
deterministic run has only one transcript, so \(h=h'\).  Hence distinct cells
are disjoint.  Also
\(\omega_h\in C_h\) for every \(h\), because a representative has the same
injection as itself and agrees with itself on every coordinate of \(R_h\).

The exact-recorded equivalence (1') follows from the same uniqueness argument.
Let \(\omega\in C_h\).  Since \(C_h\) gives the same injection as \(\omega_h\)
and the same truth table on \(R_h=R(\omega_h)\), deterministic replay
uniqueness applied to \((\omega_h,\omega)\) gives
\[
R(\omega)=R(\omega_h)=R_h.
\tag{6a}
\]
For the first inclusion, if \(e\in R_h\), then (6a) gives
\(e\in R(\omega)\), and (3) gives
\[
e\in\noderef{TwoBitePreTerminalRecordedPairs}(\omega,\varepsilon,I).
\]
For the reverse inclusion, if
\[
e\in\noderef{TwoBitePreTerminalRecordedPairs}(\omega,\varepsilon,I),
\]
then (3) gives \(e\in R(\omega)\), and (6a) gives \(e\in R_h\).  Thus, for
every tagged coordinate \(e\),
\[
e\in R_h
\Longleftrightarrow
e\in\noderef{TwoBitePreTerminalRecordedPairs}(\omega,\varepsilon,I).
\tag{6}
\]

Now prove the replay-preservation clause.  Let \(\omega\in C_h\), and let
\(\omega'\) have the same injection as \(\omega\) and the same edge values as
\(\omega\) on every coordinate of \(R_h\).  Projection containment agrees
between \(\omega\) and \(\omega'\), because
\(\noderef{TwoBiteProjectionContains}\) uses only the red and blue projections
of the common injection.  The fibers
\(\noderef{TwoBiteBaseFiber}(-,x)\cap I\) also agree, because the fibers are
defined by those same projections.  Therefore every
\(\noderef{TwoBiteStageF1}\) predicate agrees.

Fix \(x\).  If \(x\) is not projection-contained, then
\(\noderef{TwoBiteStageF2}(-,\varepsilon,I,x)\) is false for both samples.  If
\(x\) is projection-contained, the possible \(y\)'s counted by \(F_2\) are the
same in both samples because the \(F_1\) table agrees.  If a candidate has the
opposite color from \(x\), then
\(\noderef{TwoBiteBaseAdj}(-,x,y)\) is false in both samples by definition of
\(\noderef{TwoBiteBaseAdj}\).  If \(y=x\), then the same adjacency is a loop
and is false in both samples because the base red and blue graphs are simple.
For every remaining same-color candidate \(y\in F_1\) with \(y\ne x\), the
definition of \(\noderef{TwoBiteStageRevealedVertex}\) makes \(y\)
staged-revealed in \(\omega\), because its \(F_1\)-disjunct holds.  The
endpoint \(x\) is projection-contained in \(\omega\) by the present case.
Thus the incident-coordinate branch of
\(\noderef{TwoBitePreTerminalRecordedPairs}(\omega,\varepsilon,I)\) contains
\([x,y]\).  By (6), \([x,y]\in R_h\).  The assumed \(R_h\)-agreement gives the
same edge value for \([x,y]\), hence the same truth value of
\(\noderef{TwoBiteBaseAdj}(-,x,y)\).  Thus the finite \(F_2\) count is
identical for \(\omega\) and \(\omega'\), so
\[
\noderef{TwoBiteStageF2}(\omega,\varepsilon,I,x)
\Longleftrightarrow
\noderef{TwoBiteStageF2}(\omega',\varepsilon,I,x).
\]
Combining projection containment, \(F_1\), and \(F_2\) in (2), every
\(\noderef{TwoBiteStageRevealedVertex}\) predicate agrees between \(\omega\)
and \(\omega'\).

Finally, \(\noderef{TwoBiteProjectionPairTouched}\) is defined by cases on the
tagged coordinate and says that one of the two same-color endpoints is
staged-revealed.  Since staged-revealed vertices agree for all base vertices,
the touched predicate agrees for every tagged coordinate.

It remains to prove the terminal-exclusion clause.  Let
\[
e\in\noderef{TwoBiteStagedOpenPairs}(\omega,\varepsilon,I).
\]
Unfolding \(\noderef{TwoBiteStagedOpenPairs}\), the coordinate \(e\) lies in
\(\noderef{TwoBiteProjectionPairSet}(\omega,I)\), is not
\(\noderef{TwoBiteProjectionPairSameColorClosed}\), is not touched, and is not
in \(\noderef{TwoBitePreTerminalRecordedPairs}(\omega,\varepsilon,I)\).
From \(e\in\noderef{TwoBiteProjectionPairSet}(\omega,I)\), unfold the
projected-pair-set definition.  In the red case,
\[
e=\operatorname{red}(q)
\quad\text{with}\quad
q\in\noderef{TwoBiteBasePairs}(I.\operatorname{image}\pi_R^\omega),
\]
and in the blue case the analogous statement holds with \(\pi_B^\omega\).
Unfolding \(\noderef{TwoBiteBasePairs}\) and
\(\noderef{TwoBitePairsInSet}\), the pair \(q\) is represented in increasing
non-loop order.  Therefore \(e\in
\noderef{TwoBiteTerminalCoordinateUniverse}(m)\).

If \(e\in R_h\), then the exact-recorded equivalence (6), applied to
\(\omega\in C_h\), gives
\[
e\in\noderef{TwoBitePreTerminalRecordedPairs}(\omega,\varepsilon,I),
\]
contradicting the recorded-coordinate exclusion in staged-openness.  Hence
\(e\notin R_h\).  This proves the terminal-universe and unrecorded conclusions
and completes the construction of the finite disjoint pre-terminal history
partition.
\end{proof}
